\documentclass[11pt]{article}

\usepackage[margin=1in]{geometry}
\usepackage{amsmath,amssymb,amsthm}
\usepackage{algorithm}
\usepackage{algpseudocode}
\usepackage{enumitem}
\usepackage{hyperref}
\usepackage{xcolor}
\usepackage{tikz}
\usepackage{booktabs}

\hypersetup{colorlinks=true, linkcolor=blue, urlcolor=blue, citecolor=blue}

\newtheorem{theorem}{Theorem}[section]
\newtheorem{lemma}[theorem]{Lemma}
\newtheorem{proposition}[theorem]{Proposition}
\newtheorem{corollary}[theorem]{Corollary}
\newtheorem{definition}[theorem]{Definition}
\newtheorem{example}[theorem]{Example}

\title{TOAE209/TOAH209 -- Design and Analysis of Algorithms\\[4pt]
\large Week 10: Network Flow I: Maximum Flow and Minimum Cut}
\author{Foreign Trade University}
\date{}

\begin{document}
\maketitle
\tableofcontents

\section{Introduction: Flow as a Modelling Tool}

This week we study one of the most versatile algorithmic frameworks in the entire course:
\textbf{network flow}. The setting is intuitive -- imagine a network of pipes (a directed
graph), each pipe with a fixed maximum throughput (a \emph{capacity}), and ask: how much
``stuff'' (water, traffic, data, abstract units) can we route from a designated
\emph{source} to a designated \emph{sink} without exceeding any pipe's capacity? This is
the \textbf{Maximum Flow} problem.

What makes flow so important is not the pipe story but the astonishing range of problems
that turn out to be flow in disguise. Bipartite matching, edge-disjoint paths, project
selection, image segmentation, baseball elimination, and many scheduling problems all
reduce to a single max-flow computation. Two of this week's exercises (Problems 9 and 11)
develop the matching and edge-disjoint-paths reductions directly; next week is devoted
almost entirely to such applications.

Underlying everything is a single, beautiful duality result -- the \textbf{Max-Flow
Min-Cut Theorem} -- which says that the maximum amount you can push through a network
exactly equals the capacity of the cheapest ``bottleneck'' that separates source from
sink. This theorem ties together an \emph{optimization} problem (maximize flow) and a
\emph{combinatorial} problem (find the minimum cut), and the proof itself \emph{is} the
correctness proof of our main algorithm. We follow Kleinberg \& Tardos, Chapter 7.

\section{Flow Networks, Flows, and Their Value}

\subsection{Flow networks}

\begin{definition}[Flow network]
A \emph{flow network} is a directed graph $G = (V, E)$ together with:
\begin{itemize}
    \item a non-negative \emph{capacity} $c(e) \ge 0$ on each edge $e \in E$ (we assume
    integer capacities unless stated otherwise);
    \item a distinguished \emph{source} node $s \in V$ and \emph{sink} node $t \in V$,
    with $s \ne t$.
\end{itemize}
For simplicity we assume no edge enters $s$ and no edge leaves $t$, and that there are no
parallel edges in the same direction (a pair $u \to v$ and $v \to u$ is allowed).
\end{definition}

We think of $c(e)$ as the maximum rate at which material can travel along edge $e$.

\subsection{Flows}

\begin{definition}[$s$-$t$ flow]
An \emph{$s$-$t$ flow} is a function $f : E \to \mathbb{R}_{\ge 0}$ satisfying two
conditions:
\begin{enumerate}[label=(\roman*)]
    \item \textbf{Capacity constraints.} For every edge $e$, $\;0 \le f(e) \le c(e)$.
    \item \textbf{Conservation constraints.} For every node $v \ne s, t$,
    \[
        \sum_{e \text{ into } v} f(e) \;=\; \sum_{e \text{ out of } v} f(e),
    \]
    i.e.\ flow in equals flow out (nothing is created or destroyed at internal nodes).
\end{enumerate}
\end{definition}

\begin{definition}[Value of a flow]
The \emph{value} of a flow $f$, written $v(f)$, is the net flow out of the source:
\[
    v(f) \;=\; \sum_{e \text{ out of } s} f(e) \;-\; \sum_{e \text{ into } s} f(e).
\]
(With our assumption that no edge enters $s$, this is just the total flow leaving $s$.)
\end{definition}

The \textbf{Maximum-Flow Problem} is: given a flow network, find a flow of maximum value.
Problem 1 in this week's exercises asks you to write a \emph{flow feasibility checker}
that confirms the two conditions above and returns $v(f)$ -- a good way to internalize the
definitions before computing flows.

\subsection{A first conservation identity}

A flow's value can be measured at the source \emph{or} at the sink -- they always agree.

\begin{lemma}
\label{lem:value-at-sink}
For any $s$-$t$ flow $f$, the net flow out of $s$ equals the net flow into $t$.
\end{lemma}

\begin{proof}
Sum the conservation equation ``flow in $=$ flow out'' over \emph{all} nodes $v \ne s, t$,
and note that every edge $e = (u,w)$ with $u, w \notin \{s,t\}$ contributes $+f(e)$ once
(as flow out of $u$) and $-f(e)$ once (as flow into $w$), cancelling. After cancellation,
only terms touching $s$ and $t$ survive, and rearranging gives net-out$(s)=$ net-in$(t)$.
This is a special case of the more general cut identity in
Lemma~\ref{lem:flow-across-cut}.
\end{proof}

\section{Cuts and Cut Capacity}

\begin{definition}[$s$-$t$ cut]
An \emph{$s$-$t$ cut} is a partition $(S, T)$ of the node set $V$ into two sets with
$s \in S$ and $t \in T$ (so $T = V \setminus S$). The \emph{capacity} of the cut is the
total capacity of edges directed from $S$ to $T$:
\[
    c(S, T) \;=\; \sum_{\substack{e = (u,v) \\ u \in S,\; v \in T}} c(e).
\]
\end{definition}

Crucially, only edges going \emph{from $S$ to $T$} count; edges from $T$ back to $S$ are
free. Problem 2 asks you to compute $c(S,T)$ given $S$.

Cuts give an immediate \emph{upper bound} on any flow: no matter how cleverly you route
flow, every unit must cross from $S$ to $T$ at some point, and the $S \to T$ edges can
carry at most $c(S,T)$ in total.

\begin{lemma}[Flow across a cut]
\label{lem:flow-across-cut}
Let $f$ be any $s$-$t$ flow and $(S,T)$ any $s$-$t$ cut. Then
\[
    v(f) \;=\; \sum_{\substack{e \text{ out of } S}} f(e) \;-\; \sum_{\substack{e \text{ into } S}} f(e),
\]
where ``out of $S$'' means edges from $S$ to $T$ and ``into $S$'' means edges from $T$
to $S$.
\end{lemma}

\begin{proof}
By definition $v(f) = \sum_{e\text{ out of }s} f(e) - \sum_{e\text{ into }s} f(e)$. For
every node $v \in S$ with $v \ne s$, conservation gives
$\sum_{e\text{ out of }v} f(e) - \sum_{e\text{ into }v} f(e) = 0$. Add all these (one per
$v \in S$) to $v(f)$; the value is unchanged because we added zeros. Now group the sum by
edges. An edge with both endpoints in $S$ contributes $+f(e)$ (out of its tail) and
$-f(e)$ (into its head), cancelling. An edge from $S$ to $T$ contributes only $+f(e)$; an
edge from $T$ to $S$ contributes only $-f(e)$. The result is exactly the claimed
expression.
\end{proof}

\begin{corollary}[Weak duality]
\label{cor:weak-duality}
For any flow $f$ and any $s$-$t$ cut $(S,T)$, $\;v(f) \le c(S,T)$. Consequently,
\[
    \max_f v(f) \;\le\; \min_{(S,T)} c(S,T).
\]
\end{corollary}

\begin{proof}
By Lemma~\ref{lem:flow-across-cut}, $v(f) = (\text{flow }S\!\to\!T) - (\text{flow
}T\!\to\!S) \le (\text{flow }S\!\to\!T) \le \sum_{e: S\to T} c(e) = c(S,T)$, using
$f(e) \ge 0$ for the first inequality and $f(e) \le c(e)$ for the second. Taking the max
over flows on the left and the min over cuts on the right gives the second statement.
\end{proof}

The Max-Flow Min-Cut Theorem (Section~\ref{sec:mfmc}) will show this inequality is
actually an \emph{equality}.

\section{The Ford-Fulkerson Method}

\subsection{Why greedy alone fails}

A first instinct: repeatedly find a path from $s$ to $t$ with spare capacity and push as
much flow as the path's bottleneck allows. This is on the right track, but pure greedy
gets \emph{stuck}: an early path choice can ``block'' a better routing, and -- with no way
to \emph{undo} flow -- the algorithm halts below the optimum. The classic example is the
diamond network $s \to a$, $s \to b$, $a \to b$, $a \to t$, $b \to t$ (all capacity 1):
greedily routing $s \to a \to b \to t$ saturates the middle and leaves value 1, but the
true maximum is 2. The fix is the \emph{residual graph}, which lets us cancel previously
routed flow. (Problems 5 and 6 both compute the correct value 2 on this instance.)

\subsection{The residual graph}

\begin{definition}[Residual graph]
Given a flow $f$ in $G$, the \emph{residual graph} $G_f$ has the same nodes and the
following edges:
\begin{itemize}
    \item For each edge $e = (u,v)$ of $G$ with $f(e) < c(e)$, a \emph{forward} residual
    edge $(u,v)$ with \emph{residual capacity} $c(e) - f(e)$ -- the spare capacity we can
    still push.
    \item For each edge $e = (u,v)$ of $G$ with $f(e) > 0$, a \emph{backward} residual
    edge $(v,u)$ with residual capacity $f(e)$ -- representing the option to \emph{cancel}
    up to $f(e)$ units already sent along $e$.
\end{itemize}
\end{definition}

In our implementation (see \texttt{starter\_code.py}), every edge is stored together with
a paired reverse edge of capacity $0$; pushing $d$ units on an edge does \texttt{flow +=
d} on it and \texttt{flow -= d} on its reverse, so residual capacities in both directions
stay consistent automatically. Problem 3 asks you to construct $G_f$ explicitly from a
flow $f$.

\subsection{Augmenting paths}

\begin{definition}[Augmenting path]
An \emph{augmenting path} is a simple path from $s$ to $t$ in the residual graph $G_f$.
Its \emph{bottleneck} is the minimum residual capacity of its edges. \emph{Augmenting}
along it means pushing \texttt{bottleneck} units of flow on each of its edges (increasing
$f$ on forward edges, decreasing $f$ on backward ones).
\end{definition}

\begin{lemma}
\label{lem:augment-valid}
If $f$ is a flow and $P$ is an augmenting path with bottleneck $b > 0$, then augmenting
$f$ along $P$ yields a new flow $f'$ with $v(f') = v(f) + b$.
\end{lemma}

\begin{proof}[Proof sketch]
On each forward residual edge we add $b \le c(e) - f(e)$, keeping $f' \le c$; on each
backward residual edge we subtract $b \le f(e)$, keeping $f' \ge 0$. Conservation is
preserved at every internal node of $P$ because the path enters and leaves the node by
exactly one residual edge each, changing in- and out-flow by the same $b$. The first edge
of $P$ leaves $s$ (a forward edge, since no edge enters $s$), so $v$ increases by exactly
$b$.
\end{proof}

\subsection{The Ford-Fulkerson method}

\begin{algorithm}[h]
\caption{Ford-Fulkerson (generic augmenting-path method)}
\begin{algorithmic}[1]
\State Initialize $f(e) \gets 0$ for all edges $e$
\While{there exists an augmenting path $P$ in the residual graph $G_f$}
    \State $b \gets$ bottleneck residual capacity along $P$
    \State Augment $f$ along $P$ by $b$ (update forward and backward residual edges)
\EndWhile
\State \Return $f$
\end{algorithmic}
\end{algorithm}

Ford-Fulkerson is a \emph{method}, not a fully specified algorithm, because it does not
say \emph{which} augmenting path to pick. Problem 6 implements the version that finds
augmenting paths by \emph{depth-first search}; Section~\ref{sec:edmonds-karp} below pins
down the path choice to get a good running-time bound.

\section{The Max-Flow Min-Cut Theorem}
\label{sec:mfmc}

We now prove the central result. The argument is elegant: when Ford-Fulkerson stops, the
nodes reachable from $s$ in the final residual graph define a cut whose capacity equals
the flow value -- which, by weak duality, forces both to be optimal simultaneously.

\begin{theorem}[Max-Flow Min-Cut Theorem]
\label{thm:mfmc}
In every flow network, the maximum value of an $s$-$t$ flow equals the minimum capacity of
an $s$-$t$ cut. Moreover, a flow $f$ is maximum if and only if its residual graph $G_f$
contains \emph{no} augmenting path (no $s$-$t$ path).
\end{theorem}

\begin{proof}
We prove the following three statements are equivalent for a flow $f$:
\begin{enumerate}[label=(\roman*)]
    \item there is a cut $(S,T)$ with $c(S,T) = v(f)$;
    \item $f$ is a maximum flow;
    \item $G_f$ contains no augmenting path.
\end{enumerate}

\textbf{(i) $\Rightarrow$ (ii).} If $c(S,T) = v(f)$ for some cut, then by weak duality
(Corollary~\ref{cor:weak-duality}) every flow $f'$ satisfies $v(f') \le c(S,T) = v(f)$,
so $f$ is maximum (and $(S,T)$ is simultaneously a minimum cut).

\textbf{(ii) $\Rightarrow$ (iii).} Contrapositive: if $G_f$ has an augmenting path, then
by Lemma~\ref{lem:augment-valid} we can increase $v(f)$, so $f$ is not maximum.

\textbf{(iii) $\Rightarrow$ (i).} Suppose $G_f$ has no $s$-$t$ path. Let
\[
    S \;=\; \{\, v \in V : v \text{ is reachable from } s \text{ in } G_f \,\},
    \qquad T = V \setminus S.
\]
Then $s \in S$, and $t \notin S$ (else there would be an augmenting path), so $(S,T)$ is a
genuine $s$-$t$ cut. Consider any edge $e=(u,v)$ of the \emph{original} graph with
$u \in S$ and $v \in T$. We must have $f(e) = c(e)$ (the edge is \emph{saturated});
otherwise $G_f$ would contain a forward residual edge $(u,v)$, making $v$ reachable from
$s$ -- contradicting $v \in T$. Likewise, any original edge $e'=(v,u)$ with $v \in T$,
$u \in S$ must have $f(e') = 0$; otherwise $G_f$ would contain a backward residual edge
$(u,v)$, again making $v$ reachable. Now apply Lemma~\ref{lem:flow-across-cut}:
\[
    v(f) \;=\; \underbrace{\sum_{e:\, S \to T} f(e)}_{=\,\sum c(e)\,=\,c(S,T)}
            \;-\; \underbrace{\sum_{e:\, T \to S} f(e)}_{=\,0}
        \;=\; c(S,T).
\]
This proves (i). Since (i), (ii), (iii) are equivalent, the maximum flow value equals
$c(S,T)$ for this particular cut, and by weak duality no cut can be smaller -- so it is the
\emph{minimum} cut, and max-flow $=$ min-cut.
\end{proof}

The construction in (iii) $\Rightarrow$ (i) is \emph{algorithmic}: after computing a
maximum flow, do one BFS from $s$ in the final residual graph; the reached set $S$ is one
side of a minimum cut. Problem 8 implements exactly this extraction and verifies
$c(S,T) = v(f)$ on random networks.

\section{Edmonds-Karp and the Integrality Theorem}
\label{sec:edmonds-karp}

\subsection{Termination and integrality}

\begin{theorem}[Integrality Theorem]
\label{thm:integrality}
If all capacities are integers, then Ford-Fulkerson terminates, and the maximum flow it
produces is \emph{integral}: $f(e) \in \mathbb{Z}_{\ge 0}$ on every edge.
\end{theorem}

\begin{proof}
With integer capacities, the all-zero starting flow is integral. By induction, if $f$ is
integral, then every residual capacity is an integer, so each augmenting path's bottleneck
$b$ is a positive integer, and augmenting keeps $f$ integral. Each augmentation increases
$v(f)$ by $b \ge 1$, and $v(f)$ is bounded above by $c(\{s\}, V\setminus\{s\})$ (the
capacity of the source cut), so there can be only finitely many augmentations. Hence the
method terminates, with an integral maximum flow.
\end{proof}

Integrality is not just a curiosity -- it is precisely what makes the matching and
edge-disjoint-paths reductions (Problems 9 and 11) work: a unit-capacity flow comes back
as a $0/1$ assignment that we can read off as ``edge chosen'' or ``edge not chosen.''
Problem 12 confirms integrality empirically: on integer-capacity networks, the computed
flow is an integer on every edge.

\subsection{The Edmonds-Karp BFS rule}

Generic Ford-Fulkerson can be slow (even non-terminating with irrational capacities, and
exponential with adversarial integer choices). \textbf{Edmonds and Karp} fix the
path-selection rule: \emph{always augment along a shortest augmenting path}, i.e.\ one with
the fewest edges, found by BFS.

\begin{algorithm}[h]
\caption{Edmonds-Karp (BFS augmenting paths)}
\begin{algorithmic}[1]
\State $f(e) \gets 0$ for all $e$
\While{BFS from $s$ in $G_f$ reaches $t$}
    \State Let $P$ be the BFS shortest $s$-$t$ path; $b \gets$ its bottleneck
    \State Augment $f$ along $P$ by $b$
\EndWhile
\State \Return $f$
\end{algorithmic}
\end{algorithm}

Problem 4 implements the BFS that returns one such shortest augmenting path (or
\texttt{None} when $t$ is unreachable); Problem 5 wraps it into the full Edmonds-Karp
algorithm.

\begin{theorem}[Edmonds-Karp bound]
\label{thm:ek-bound}
With the shortest-augmenting-path rule, Ford-Fulkerson performs $O(VE)$ augmentations,
each costing $O(E)$ time for the BFS, for a total running time of $O(VE^2)$ -- independent
of the capacity values.
\end{theorem}

\begin{proof}[Proof idea]
Two facts drive the bound. (1) Under BFS augmentation, the BFS distance
$\mathrm{dist}_{G_f}(s, v)$ from $s$ to any node $v$ \emph{never decreases} over the course
of the algorithm. (2) An edge $(u,v)$ becomes \emph{saturated} (the bottleneck of the
chosen path) only when it lies on a shortest path, i.e.\ $\mathrm{dist}(s,v) =
\mathrm{dist}(s,u)+1$; before it can be saturated again, flow must be pushed back along
$(v,u)$, which requires $\mathrm{dist}(s,u)$ to have grown by at least $2$. Since distances
range in $\{0, 1, \ldots, V-1, \infty\}$, each edge can be saturated at most $O(V)$ times,
giving $O(VE)$ total saturations, hence $O(VE)$ augmentations (each augmentation saturates
at least one edge). Each BFS is $O(V+E) = O(E)$, so the total is $O(VE^2)$. The full
distance-monotonicity argument is in Kleinberg \& Tardos, Section 7.3.
\end{proof}

\section{Putting the Pieces Together: Verification}

A theme of this course is \emph{verifying} that an algorithm's output is actually
correct, not just trusting it. Network flow is especially well-suited to this:
\begin{itemize}
    \item \textbf{Cross-checking algorithms.} Since the max-flow value is \emph{unique}
    (it equals the min-cut), two different correct methods must agree. Problem 7 runs
    Edmonds-Karp (BFS) and DFS Ford-Fulkerson on the same random networks and confirms
    they always return the same value.
    \item \textbf{Certificate of optimality.} A flow is provably maximum if you can
    exhibit a cut of equal capacity (Theorem~\ref{thm:mfmc}). Problem 8 extracts that cut
    and checks $c(S,T) = v(f)$ -- a self-certifying proof of optimality.
    \item \textbf{Structural invariants.} Problem 1 checks the defining flow conditions;
    Problem 12 checks integrality. These act as assertions that any future change to the
    code must preserve.
\end{itemize}

\section{Summary and Looking Ahead}

This week introduced the \textbf{network flow} framework and its central duality:

\begin{itemize}
    \item A \textbf{flow network} is a directed graph with edge capacities, a source $s$,
    and a sink $t$; a \textbf{flow} respects capacities and conserves material at internal
    nodes; its \textbf{value} is the net flow out of $s$.
    \item An \textbf{$s$-$t$ cut} $(S,T)$ has capacity equal to the total capacity of
    $S \to T$ edges, and \emph{every} cut upper-bounds \emph{every} flow (weak duality).
    \item \textbf{Ford-Fulkerson} repeatedly augments along $s$-$t$ paths in the
    \textbf{residual graph}, whose backward edges allow flow to be cancelled.
    \item The \textbf{Max-Flow Min-Cut Theorem} states that the maximum flow value equals
    the minimum cut capacity; the proof simultaneously yields an algorithm for extracting
    a minimum cut from any maximum flow.
    \item \textbf{Edmonds-Karp} (augment along shortest BFS paths) runs in $O(VE^2)$, and
    the \textbf{Integrality Theorem} guarantees integer flows for integer capacities.
\end{itemize}

We also saw how to \emph{verify} flows and cuts, and previewed two reductions -- bipartite
matching (Problem 9) and edge-disjoint paths via Menger's theorem (Problem 11), plus the
super-source/super-sink trick for multiple sources and sinks (Problem 10).

Next week (Network Flow II/III, Kleinberg \& Tardos Ch.\ 7 continued) we turn the crank in
earnest on \textbf{applications of flow}: maximum bipartite matching and Hall's theorem,
disjoint paths and Menger's theorem in full, and a gallery of modelling problems
(survey design, project selection / image segmentation, baseball elimination) that all
collapse to a single max-flow or min-cut computation.

\end{document}
